\documentclass[
    language=english,
    doctype=manual,
    institution=none,
    titlestyle=book,
    oneside,
]{../../omnilatex}

\RequirePackage{config/document-settings}

\addbibresource{../../bib/bibliography.bib}

\begin{document}

\maketitle

\tableofcontents

%=============================================================================
\chapter{Introduction}
\label{sec:introduction}
%=============================================================================

This guide helps you diagnose and resolve common DataPipe issues. Each
section includes the error description, diagnostic steps, root cause
analysis, and resolution. Use the error code table to quickly locate the
relevant section.

\section{How to Use This Guide}

\begin{enumerate}
    \item Identify the error code or symptom from the logs.
    \item Look up the error code in the table below.
    \item Follow the diagnostic flowchart for your issue category.
    \item Apply the recommended resolution.
    \item If the issue persists, contact support with the collected
        diagnostic information.
\end{enumerate}

%=============================================================================
\chapter{Error Code Reference}
\label{sec:error-codes}
%=============================================================================

\begin{table}[htbp]
    \centering
    \caption{DataPipe error codes and descriptions.}
    \label{tab:error-codes}
    \begin{tabular}{@{}p{2cm}p{4cm}p{2cm}p{3cm}@{}}
        \toprule
        Code & Description & Severity & Section \\
        \midrule
        E001 & Kafka connection refused & Critical & \cref{sec:e001} \\
        E002 & Pipeline configuration invalid & High & \cref{sec:e002} \\
        E003 & Sink write failure & Critical & \cref{sec:e003} \\
        E004 & Transform execution timeout & Medium & \cref{sec:e004} \\
        E005 & Consumer lag threshold exceeded & High & \cref{sec:e005} \\
        E006 & Out of memory & Critical & \cref{sec:e006} \\
        E007 & Schema evolution conflict & High & \cref{sec:e007} \\
        E008 & Authentication failure & Critical & \cref{sec:e008} \\
        E009 & Disk space exhausted & Critical & \cref{sec:e009} \\
        E010 & Checkpoint corruption & High & \cref{sec:e010} \\
        \bottomrule
    \end{tabular}
\end{table}

%=============================================================================
\chapter{Diagnostic Flowcharts}
\label{sec:flowcharts}
%=============================================================================

\section{Pipeline Not Starting}

\begin{figure}[htbp]
    \centering
    \begin{tikzpicture}[
        node distance=1cm and 1.5cm,
        decision/.style={diamond, draw, text width=2.5cm,
                         text centered, inner sep=1pt,
                         fill=yellow!20},
        action/.style={rectangle, draw, rounded corners,
                       minimum width=3cm, minimum height=0.8cm,
                       text centered, font=\small, fill=blue!15},
        result/.style={rectangle, draw, rounded corners,
                       minimum width=3cm, minimum height=0.8cm,
                       text centered, font=\small,
                       fill=green!20, thick},
        arr/.style={->, >=stealth, thick},
        label/.style={font=\scriptsize},
    ]
        \node[action] (start) {Pipeline won't start};
        \node[decision, below=of start] (config) {Config valid?};
        \node[action, right=of config] (fix-config) {Fix configuration};
        \node[decision, below=of config] (kafka) {Kafka reachable?};
        \node[action, right=of kafka] (fix-kafka) {Check network};
        \node[decision, below=of kafka] (perm) {Permissions OK?};
        \node[action, right=of perm] (fix-perm) {Fix credentials};
        \node[decision, below=of perm] (resources) {Resources OK?};
        \node[action, right=of resources] (fix-res) {Scale up};
        \node[result, below=of resources] (ok) {Pipeline starts};

        \draw[arr] (start) -- (config);
        \draw[arr] (config) -- node[label, right] {No} (fix-config);
        \draw[arr] (config) -- node[label, below] {Yes} (kafka);
        \draw[arr] (kafka) -- node[label, right] {No} (fix-kafka);
        \draw[arr] (kafka) -- node[label, below] {Yes} (perm);
        \draw[arr] (perm) -- node[label, right] {No} (fix-perm);
        \draw[arr] (perm) -- node[label, below] {Yes} (resources);
        \draw[arr] (resources) -- node[label, right] {No} (fix-res);
        \draw[arr] (resources) -- node[label, below] {Yes} (ok);
    \end{tikzpicture}
    \caption{Diagnostic flowchart: pipeline not starting.}
    \label{fig:flow-pipeline-start}
\end{figure}

\section{High Consumer Lag}

\begin{figure}[htbp]
    \centering
    \begin{tikzpicture}[
        node distance=1cm and 1.5cm,
        decision/.style={diamond, draw, text width=2.5cm,
                         text centered, inner sep=1pt,
                         fill=yellow!20},
        action/.style={rectangle, draw, rounded corners,
                       minimum width=3cm, minimum height=0.8cm,
                       text centered, font=\small, fill=blue!15},
        result/.style={rectangle, draw, rounded corners,
                       minimum width=3cm, minimum height=0.8cm,
                       text centered, font=\small,
                       fill=green!20, thick},
        arr/.style={->, >=stealth, thick},
        label/.style={font=\scriptsize},
    ]
        \node[action] (start) {Consumer lag detected};
        \node[decision, below=of start] (throughput)
            {Source throughput normal?};
        \node[action, right=of throughput] (source-issue)
            {Investigate source};
        \node[decision, below=of throughput] (cpu)
            {CPU $> 80\%$?};
        \node[action, right=of cpu] (scale) {Scale consumers};
        \node[decision, below=of cpu] (slow-transform)
            {Transform slow?};
        \node[action, right=of slow-transform] (optimise)
            {Optimise transform};
        \node[decision, below=of slow-transform] (sink)
            {Sink bottleneck?};
        \node[action, right=of sink] (fix-sink)
            {Optimise sink writes};
        \node[result, below=of sink] (ok) {Lag recovered};

        \draw[arr] (start) -- (throughput);
        \draw[arr] (throughput) -- node[label, right] {No} (source-issue);
        \draw[arr] (throughput) -- node[label, below] {Yes} (cpu);
        \draw[arr] (cpu) -- node[label, right] {Yes} (scale);
        \draw[arr] (cpu) -- node[label, below] {No} (slow-transform);
        \draw[arr] (slow-transform) -- node[label, right] {Yes} (optimise);
        \draw[arr] (slow-transform) -- node[label, below] {No} (sink);
        \draw[arr] (sink) -- node[label, right] {Yes} (fix-sink);
        \draw[arr] (sink) -- node[label, below] {No} (ok);
    \end{tikzpicture}
    \caption{Diagnostic flowchart: high consumer lag.}
    \label{fig:flow-consumer-lag}
\end{figure}

%=============================================================================
\chapter{Common Issues and Solutions}
\label{sec:common-issues}
%=============================================================================

\section{E001: Kafka Connection Refused}
\label{sec:e001}

\textbf{Symptom:} Pipeline fails to start with error
\texttt{E001: Kafka connection refused}.

\textbf{Diagnostic steps:}

\begin{lstlisting}[language=bash]
# Verify Kafka is running
kubectl get pods -l app=kafka -n datapipe

# Test connectivity
kubectl exec datapipe-0 -n datapipe -- \\
    kafka-broker-api-versions.sh \\
    --bootstrap-server kafka:9092

# Check DNS resolution
kubectl exec datapipe-0 -n datapipe -- nslookup kafka
\end{lstlisting}

\textbf{Common causes and fixes:}

\begin{description}
    \item[Kafka pods not running] Restart Kafka pods or check
        Kafka operator logs.
    \item[DNS misconfiguration] Verify the Kafka service name in the
        pipeline configuration matches the Kubernetes service.
    \item[Network policy blocking] Check Kubernetes NetworkPolicy
        resources in the \texttt{datapipe} namespace.
    \item[Firewall rules] Verify security group rules allow traffic
        on port 9092 between application and Kafka nodes.
\end{description}

\section{E002: Pipeline Configuration Invalid}
\label{sec:e002}

\textbf{Symptom:} Pipeline fails to start with error
\texttt{E002: Pipeline configuration invalid}.

\textbf{Diagnostic steps:}

\begin{lstlisting}[language=bash]
# Validate the pipeline configuration
dp pipeline validate <pipeline-name> --verbose

# Check for YAML syntax errors
yamllint <config-file>.yaml
\end{lstlisting}

\textbf{Common causes:}

\begin{itemize}
    \item Missing required fields in the pipeline definition.
    \item YAML indentation errors.
    \item Invalid connector type name.
    \item Mismatched environment variable references.
\end{itemize}

\section{E003: Sink Write Failure}
\label{sec:e003}

\textbf{Symptom:} Pipeline reports errors writing to the sink, records
are retried and eventually dropped.

\textbf{Diagnostic steps:}

\begin{lstlisting}[language=bash]
# Check sink connectivity
kubectl exec datapipe-0 -n datapipe -- \\
    psql -h postgresql -U datapipe -c "SELECT 1;"

# Check PostgreSQL replication status
kubectl exec postgresql-primary-0 -n datapipe -- \\
    psql -U postgres -c "SELECT * FROM pg_stat_replication;"

# Check disk space on PostgreSQL
kubectl exec postgresql-primary-0 -n datapipe -- \\
    df -h /var/lib/postgresql/data
\end{lstlisting}

\textbf{Common causes and fixes:}

\begin{description}
    \item[Table does not exist] Run the schema migration:
        \texttt{dp db migrate run}.
    \item[Disk full] Expand the persistent volume or archive old data.
    \item[Connection pool exhausted] Increase
        \texttt{sink.max\_connections} in the pipeline config.
    \item[Permission denied] Grant \texttt{INSERT} permission on the
        target table.
\end{description}

\section{E006: Out of Memory}
\label{sec:e006}

\textbf{Symptom:} DataPipe process crashes with
\texttt{java.lang.OutOfMemoryError}.

\textbf{Diagnostic steps:}

\begin{lstlisting}[language=bash]
# Check pod memory usage
kubectl top pods -n datapipe

# Check JVM heap settings
kubectl exec datapipe-0 -n datapipe -- \\
    jcmd 1 VM.flags | grep HeapSize

# Review heap dump if available
ls -la /tmp/datapipe-heap-*.hprof
\end{lstlisting}

\textbf{Resolution:}

\begin{enumerate}
    \item Increase JVM heap via \texttt{DP\_JVM\_HEAP} environment
        variable (e.g.\ \texttt{16g}).
    \item Increase pod memory limit in the Kubernetes deployment.
    \item Profile the pipeline to identify memory-intensive transforms.
    \item Consider implementing backpressure to limit in-flight records.
\end{enumerate}

\section{E009: Disk Space Exhausted}
\label{sec:e009}

\textbf{Symptom:} DataPipe logs report \texttt{No space left on device}.

\textbf{Diagnostic steps:}

\begin{lstlisting}[language=bash]
# Check disk usage on all nodes
kubectl exec datapipe-0 -n datapipe -- df -h

# Check Kafka log retention
kubectl exec kafka-0 -n datapipe -- \\
    kafka-configs.sh --describe \\
    --bootstrap-server localhost:9092 \\
    --entity-type topics --entity-name user-events

# Find large files
kubectl exec datapipe-0 -n datapipe -- \\
    du -sh /data/datapipe/*
\end{lstlisting}

\textbf{Resolution:}

\begin{enumerate}
    \item Reduce Kafka log retention period.
    \item Archive or delete old checkpoint files.
    \item Expand the persistent volume.
    \item Enable compression for pipeline stages.
\end{enumerate}

%=============================================================================
\chapter{Collecting Diagnostic Information}
\label{sec:diagnostics}
%=============================================================================

When contacting DataPipe support, provide the following:

\begin{lstlisting}[language=bash, caption={Generate diagnostic bundle.}]
# Generate a diagnostic bundle
dp diagnostic collect --output /tmp/datapipe-diag.zip

# The bundle includes:
# - Pipeline configurations
# - Recent logs (last 24h)
# - JVM thread dumps
# - System metrics snapshot
# - Kubernetes pod descriptions
\end{lstlisting}

\begin{lstlisting}[language=bash, caption={Useful diagnostic commands.}]
# System information
uname -a
java -version
docker --version

# Kubernetes information
kubectl get events -n datapipe --sort-by=.lastTimestamp
kubectl describe pods -n datapipe

# Network diagnostics
kubectl exec datapipe-0 -n datapipe -- netstat -tlnp
kubectl exec datapipe-0 -n datapipe -- curl -s localhost:8080/health
\end{lstlisting}

\printbibliography[heading=bibintoc, title={References}]

\end{document}
